% Appendix — Length Confound Analysis
% Data source: 8-generator LOGO corpus (37,058 samples)

Surface diversity features (TTR, hapax ratio) mechanically depend on text length:
shorter texts yield higher type-token ratios by construction.
If generation depth systematically shortens text, observed depth trends could be length artifacts rather than genuine distributional shifts.
We investigate with three complementary analyses.

\paragraph{(a) Length--depth correlation.}
Spearman rank correlation between word count and generation depth is weak across all eight generators ($|\rho| < 0.14$); depth does not systematically shorten or lengthen text in our corpus.

\paragraph{(b) Length--feature correlations.}
While length--depth correlation is weak, individual features show moderate length dependence:
tail\_decay ($|\rho| = 0.38$), hapax\_ratio ($|\rho| = 0.30$), and TTR ($|\rho| = 0.29$).
Reference-model features show weaker length correlations (entropy\_mean $|\rho| = 0.09$, entropy\_std $|\rho| = 0.18$).
This motivates the length-matched evaluation below: if length mediates some depth--feature association, controlling for length will reduce classification performance.

\paragraph{(c) Length-matched classification.}
We subsample the dataset so that word-count distributions are equalized across depth classes (gen-0\,/\,gen-1\,/\,gen-2+), then retrain the binary XGBoost classifier under pooled LOGO (8-fold).
Length-matched binary F1\,=\,0.724, down from 0.770 on the unmatched set ($\Delta$\,=\,$-$0.046).
The matched subset is smaller than the full corpus; a random-subsample control at equal $N$ (below) isolates the deconfounding effect from the sample-size reduction.
To distinguish this from a sample-size effect, we run a random-subsample control at equal $N$ without length matching: F1\,=\,0.767, which confirms that the 0.046 drop reflects genuine length deconfounding rather than reduced data.

\paragraph{Conclusion.}
Length accounts for approximately 6\% of discriminative performance ($\Delta$\,=\,0.046 out of 0.770).
The majority of the depth signal is length-independent: length-matched F1 (0.724) remains far above chance and substantially above all individual detection baselines (\S\ref{sec:results-baselines}).
Surface features (TTR, hapax\_ratio) partially co-vary with length, but their depth-discriminative contribution persists after length control, consistent with generation-induced vocabulary compression operating independently of text truncation.
